\chapter{基本程式架構}
\label{ch:basicprogramstructure}

第一章已經交代程式骨架、\texttt{gcc} 編譯流程，以及 \texttt{Hello,
world!} 怎麼跑起來。 這一章不再重講 \texttt{\#include} 與前置處理／
編譯／連結，直接從輸出與輸入的實務技巧接下去：資料怎麼透過
\texttt{printf} 和 \texttt{scanf} 在程式與外界之間來回傳遞。

這一章依四個階段推進：先把 \texttt{printf} 與換行講透，再一口氣搞懂變數
的宣告與賦值，接著用 \texttt{scanf} 與回音程式確認輸入輸出的管子是通的，
最後才處理多個數字的讀寫、補零格式與 \texttt{double}。 算式本身（加減
乘除、整數除法怎麼截斷、什麼時候該用 \texttt{double} 做計算）留到下一
章運算再集中處理；這裡的目標是讓讀者能獨立寫出讀一些整數或浮點數、依
題目要求印出結果的程式骨架。

\section{用 printf 輸出文字}
\label{sec:basic-printf}

要印出文字，我們在 \texttt{printf} 的格式字串裡寫想輸出的內容；一行一
個 \texttt{printf}，就會一行一行印到螢幕上。

\begin{code}{c}{連續印出多行文字}
printf("first line\n");
printf("second line\n");
printf("third line\n");
\end{code}

格式字串用雙引號包起來；引號裡的內容會原樣輸出（除了 \texttt{\textbackslash n}
這類跳脫字元）。 如果題目要求印出固定的英文句子、而且順序與換行都已
經規定好，通常不需要變數，也不需要 \texttt{scanf}。 直接把句子寫進
\texttt{printf} 即可。

範例 \thechapter.2 依序印出三句英文；每一句各自一行，中間沒有多餘的
空白行。

\begin{example}{c}{(three-lines.c) 印出三句固定的英文}
#include <stdio.h>

int main(void) {
    printf("This is a book.\n");
    printf("It teaches C programming.\n");
    printf("Read it carefully.\n");
}
\end{example}

\iosample{}{
This is a book. \\
It teaches C programming. \\
Read it carefully. \\
}

這種純輸出、無輸入的題目，檢查的重點往往是：大小寫、標點、以及每
一行末尾有沒有照規定換行。 寫完之後用 \texttt{diff} 跟預期輸出比對，
比用肉眼掃過更可靠。

\texttt{printf} 印完一段文字之後，游標停在那段文字的右邊。 若立刻再
印下一段，兩段會黏成一長串，中間沒有空白、也沒有換行。 範例
\thechapter.3 連續印 \texttt{"17"} 和 \texttt{"29"}，螢幕上看到的是
\texttt{1729}。 這只是兩次輸出拼在一起，跟數學上的 $1729$ 無關。

\begin{example}{c}{(print-glued.c) 兩次 printf 中間沒有換行}
#include <stdio.h>

int main(void) {
    printf("17");
    printf("29");
    return 0;
}
\end{example}

\iosample{}{
1729 \\
}

補上 \texttt{\textbackslash n} 之後，每次印完都會把游標移到下一行，
如範例 \thechapter.4。 這跟上一章印 \texttt{Hello, world!} 是同一
件事。

\begin{example}{c}{(print-two-lines.c) 每段文字印完後換行}
#include <stdio.h>

int main(void) {
    printf("17\n");
    printf("29\n");
    return 0;
}
\end{example}

\iosample{}{
17 \\
29 \\
}

若希望兩段文字出現在同一行、中間用空白隔開，格式字串應寫成
\texttt{"17 29\textbackslash n"}，空白是你自己放進字串裡的字元，編
譯器不會自動插入。 題目範例輸出有幾個空白、最後有沒有換行，都要對
著抄，不要看起來差不多就過關。

\stylepoint{固定字串輸出題不需要宣告變數；把要求的文字原樣放進
\texttt{printf}，並確認 \texttt{\textbackslash n} 的位置跟題目一致。
需要換行就自己寫 \texttt{\textbackslash n}；螢幕上文字黏在一起，多半
是上一次 \texttt{printf} 沒把游標移走。}

\begin{drills}[label=drills:basic-printf]
\item 連續執行 \texttt{printf("17");} 和 \texttt{printf("29");}（中間
沒有 \texttt{\textbackslash n}），螢幕上會看到什麼？
\item 要讓 \texttt{"17"} 和 \texttt{"29"} 印在同一行、中間隔一個空
白，格式字串該怎麼寫？
\item 範例 \thechapter.4（\texttt{print-two-lines.c}）印出兩行
\texttt{17} 與 \texttt{29}，靠的是哪個字元？
\end{drills}

\section{變數與 int 型別}
\label{sec:basic-int-vars}

一旦題目說讀入一個整數，我們就需要地方暫存讀進來的值。
\emph{變數}（variable）就是程式執行期間在記憶體裡保留的一格（或多格）
空間，我們用名字（例如 \texttt{count}、\texttt{year}）存取它，型別
（例如 \texttt{int}）決定這一格有多大、能存什麼範圍的整數。

\begin{code}{c}{宣告整數變數}
int x;
\end{code}

\texttt{int x;} 向編譯器申請一個名為 \texttt{x} 的整數變數。 在
\texttt{main} 裡宣告的變數，生命週期只到 \texttt{main} 結束為止。 它
們通常配置在\emph{堆疊}（stack）上：函式被呼叫時，區域變數依序在堆
疊上開闢空間；函式返回時，這些空間一併回收。 現階段只要記住：變數名
是給人讀的，編譯器真正關心的是在記憶體的哪裡、佔幾個位元組。

每個變數可以想成同時帶著四件事：\emph{型別}（這一格能放什麼）、
\emph{名字}（我們怎麼稱呼它）、\emph{值}（這一格目前存的內容）、以
及\emph{位址}（它在記憶體裡的編號）。 \texttt{scanf} 需要的是位址，
所以才會寫 \texttt{\&x}；\texttt{printf("\%d", x)} 需要的是值，所以
寫 \texttt{x} 就夠。 把這兩種用途混在一起，是這一章後半最常見的錯。

名字不是隨便取的。 C 的\emph{識別字}（identifier）只能由英文字母、
數字和底線組成，而且不能以數字開頭；大小寫視為不同，所以
\texttt{sum} 和 \texttt{Sum} 是兩個變數。 有些單字被語言保留當關鍵
字（\texttt{int}、\texttt{return}、\texttt{if}），不能拿來當變數
名。 取名時優先用看得懂的字：\texttt{score}、\texttt{year} 比
\texttt{a1}、\texttt{tmp2} 更容易在除錯時對上號。 C 程式常見兩種風
格：小寫加底線 \texttt{student\_score}，或駝峰式 \texttt{studentScore}；
兩種都可以，重點是同一支程式裡保持一致，不要讓 \texttt{a}、
\texttt{b}、\texttt{c} 滿天飛。

\begin{center}
\begin{tikzpicture}
    \node[cell] (v) {?};
    \node[idxlbl, left=3mm of v] (name) {\texttt{x}};
    \node[idxlbl, right=3mm of v, align=left] (note) {尚未賦值};
\end{tikzpicture}
\diagcaption{變數 \texttt{x} 在概念上可以想成記憶體裡的一格；宣告後、
\texttt{scanf} 或賦值之前，內容是未定義的，不該拿來做計算。}
\label{fig:ch02-variable-box}
\end{center}

\keypoint{\texttt{int} 用來存整數；變數必須先宣告型別和名字，才能被
\texttt{scanf} 寫入或在運算式中使用。}

\stylepoint{變數名要讓人看得懂用途；\texttt{student\_score} 或
\texttt{studentScore} 都比 \texttt{a} 好維護。同一支程式裡選一種命名
風格並保持一致。}

\begin{drills}[label=drills:basic-int-vars]
\item 一個變數同時帶著哪四件事？ \texttt{scanf} 需要用到其中哪一個？
\item \texttt{int x;} 宣告完、還沒賦值之前，\texttt{x} 裡面是什麼？ 這
時候可以直接拿去 \texttt{printf} 嗎？
\item \texttt{Sum} 和 \texttt{sum} 是不是同一個變數？ 為什麼？
\end{drills}

\section{初始化與賦值}
\label{sec:basic-init-assign}

宣告變數之後，我們得把值寫進去，程式才有明確的狀態可以運作。 一個好
習慣是在宣告的當下就指定初值，讓讀程式的人一眼看出起始狀態，也避免
沒賦值就拿來計算的模糊情況。 在較複雜的程式裡，每個變數一開始都有
確定的值，除錯時也比較少追蹤到莫名其妙的殘留內容。

\begin{code}{c}{宣告並初始化}
int count = 0;
int year = 2024;
\end{code}

這段程式一邊宣告變數 \texttt{count}，一邊把它設為 $0$；\texttt{year}
則一開始就有 $2024$ 這個明確的值。 等號在這裡代表賦值：把右邊的值
寫進左邊的變數，可以想成一支向左的箭頭。 \texttt{int n = 42;} 就是
在宣告的當下就給初值；也可以先宣告、稍後再用賦值敘述改它的內容。

\begin{code}{c}{先宣告，再把值寫進變數}
int n;
n = 42;
\end{code}

如果有幾個用途相近的變數（例如同樣是迴圈用的索引、或同一組座標），
可以在同一行內一起宣告，只要型別一致即可：

\begin{code}{c}{一行宣告多個變數}
int i, j, k;
\end{code}

也可以在宣告時一併給初值：

\begin{code}{c}{宣告並初始化多個變數}
int a = 1, b = 2, c = 3;
\end{code}

範例 \thechapter.10 把宣告與初值合成一個完整程式：宣告 \texttt{answer}、
立刻給初值、再用 \texttt{\%d} 印出來。 \texttt{\%d} 的 \texttt{d} 表
示十進位（decimal）；後面那個參數必須是整數值，不是位址。

\begin{example}{c}{(show-answer.c) 設定整數並印出}
#include <stdio.h>

int main(void) {
    int answer = 42;
    printf("%d\n", answer);
    return 0;
}
\end{example}

\iosample{}{
42 \\
}

同一個意思也可以拆成先空宣告、再賦值，如範例 \thechapter.11。 執行
結果一樣是 $42$。 兩種寫法編譯器都收；初值在宣告時就寫下，比較
不容易忘記。 若宣告了卻既不初始化、也不賦值，就拿去 \texttt{printf}，
那一格裡可能是殘留的任意位元；後面會看到編譯器對這件事的警告。

\begin{example}{c}{(assign-later.c) 宣告與賦值分成兩步}
#include <stdio.h>

int main(void) {
    int n;
    n = 42;
    printf("%d\n", n);
    return 0;
}
\end{example}

\iosample{}{
42 \\
}

在這台機器上，\texttt{sizeof(int)} 的結果是 $4$，也就是一個
\texttt{int} 佔四個位元組（三十二個位元）。 這不是語言保證的唯一可
能；別的平臺上 \texttt{int} 可能更寬或更窄。 現階段只要有這個數量級
的感覺：它不是無限大的整數，太大的數字之後會在運算那一章碰到溢
位。

\keypoint{賦值是把右邊的值寫進左邊的變數。 宣告時給初值，或宣告
後再賦值，都可以；拿尚未寫入的變數去做計算，結果沒有保證。}

\stylepoint{相似用途的變數可以放在一起宣告，增加可讀性。寫作時盡量讓
名字和分組呈現變數之間的邏輯關係。}

\begin{drills}[label=drills:basic-init-assign]
\item \texttt{int n; n = 42;} 跟 \texttt{int n = 42;}，兩者最後
\texttt{n} 的值一樣嗎？
\item 宣告了 \texttt{int x;} 卻沒有初始化、也沒有賦值，就直接拿去
\texttt{printf}，這樣安全嗎？
\item 把 \texttt{int a, b, c;} 改成一行內就給三個初值，該怎麼寫？
\end{drills}

\section{用 scanf 讀入整數}
\label{sec:basic-scanf}

C 語言用 \texttt{scanf} 從標準輸入讀資料。 讀一個整數的典型寫法是：

\begin{code}{c}{讀入一個 int 並存進變數 x}
int x;
scanf("%d", &x);
\end{code}

格式字串 \texttt{"\%d"} 表示期待一個十進位整數；第二個參數必須是
變數的\emph{位址}，所以變數名前面要加 \texttt{\&}（取址運算子）。 可以
把 \texttt{\&x} 理解成請把讀到的值寫進 \texttt{x} 在記憶體裡的那一格。
\texttt{scanf} 需要知道寫到哪裡，光給變數名而不給位址是不夠的。

上一節學 \texttt{printf} 時，我們在格式字串裡用 \texttt{\textbackslash n}
控制換行。 讀入時不要照抄這個習慣：\texttt{scanf("\%d\textbackslash n",
\&n)} 是常見的地雷。 格式字串裡的空白字元（包含 \texttt{\textbackslash n}）
會一直跳過輸入裡的空白，直到看見下一個非空白字元為止；若你在讀完整數
後多寫了 \texttt{\textbackslash n}，程式可能看起來像卡住，等你在鍵盤
上多按一次才繼續。 讀整數用 \texttt{"\%d"} 就夠；換行留給 \texttt{printf}
去處理。

\keypoint{\texttt{scanf} 讀 \texttt{int} 時，格式用 \texttt{\%d}，參
數必須是 \texttt{\&變數名}；漏寫 \texttt{\&} 是初學者最常見、也最危險
的錯誤之一，可能導致程式寫壞記憶體。格式字串裡不要加 \texttt{\textbackslash n}。}

\begin{drills}[label=drills:basic-scanf]
\item 讀一個 \texttt{int} 變數 \texttt{n}，\texttt{scanf} 的格式字串
該寫成什麼？參數該怎麼寫？
\item \texttt{scanf("\%d\textbackslash n", \&n)} 這樣寫，可能會讓程式
看起來像卡住，為什麼？
\item \texttt{scanf} 的第二個參數為什麼一定要是變數的位址（加
\texttt{\&}），不能直接寫變數名？
\end{drills}

\section{回音程式：讀進來再原樣印出}
\label{sec:basic-echo}

真正解題時，數字來自標準輸入。 最小的完整迴路是：宣告、\texttt{scanf}
寫入、\texttt{printf} 讀出。 範例 \thechapter.13 什麼都不算，只把讀
到的整數印回去。 看起來無聊，卻最能確認輸入輸出這條管子是通的。 若
連回音都印錯，問題在格式字串或 \texttt{\&}，不必先懷疑後面的算式。

\begin{example}{c}{(echo-int.c) 讀入一個整數再印出}
#include <stdio.h>

int main(void) {
    int n;
    scanf("%d", &n);
    printf("%d\n", n);
    return 0;
}
\end{example}

\iosample{
100 \\
}{
100 \\
}

執行順序必須是：先有變數、再讀入、最後印出。 顛倒成先 \texttt{printf}
再 \texttt{scanf}，印出來的會是尚未寫入的內容；先 \texttt{scanf} 卻
漏掉 \texttt{\&}，則是請 \texttt{scanf} 把整數寫到一個不像位址的地
方。

管子通了之後，就可以接上一點簡單的計算。 範例 \thechapter.14 讀入兩
個整數，把它們相加後印出。 這是讀入 $\to$ 計算 $\to$ 輸出三段式程
式的最小完整版本：中間的加法本身很簡單，重點在於 \texttt{scanf} 與
\texttt{printf} 的配對是否正確；算式的細節留到下一章再展開。

\begin{example}{c}{(sum-two.c) 讀入兩個整數並輸出總和}
#include <stdio.h>

int main(void) {
    int a, b;
    scanf("%d %d", &a, &b);

    printf("%d\n", a + b);
}
\end{example}

\iosample{
3 7 \\
}{
10 \\
}

以這份輸入為例：\texttt{scanf} 依空白分隔讀到 \texttt{a = 3}、
\texttt{b = 7}，\texttt{printf} 印出 \texttt{a + b = 10}。 格式字串
\texttt{"\%d\textbackslash n"} 裡的 \texttt{\%d} 會被替換成整數值；
我們只放一個 \texttt{\%d}，所以輸出只有數字本身加換行。

我們在這台機器上用 \texttt{gcc -Wall} 編譯過一份漏寫 \texttt{\&}
的程式：\texttt{scanf("\%d", n)}。 編譯器提出兩則警告，大意是格式
期待 \texttt{int *}（指向整數的位址），實際卻收到 \texttt{int}，而
且 \texttt{n} 在呼叫當下還未初始化。 警告沒有阻止產生執行檔，但這
種程式執行期可能寫壞別的記憶體，屬於上一章說的執行期錯誤的近親。
\textbf{看到格式與型別不符的警告，先改到警告消失，再談答案對不對。}

同一個編譯器對宣告了 \texttt{int n;} 卻直接 \texttt{printf("\%d",
n)} 也會警告 \texttt{variable 'n' is uninitialized when used
here}，並建議寫 \texttt{= 0}。 初始化成 $0$ 能讓警告消失，但不表示
$0$ 就是題目要的值。 它只是一個確定的值。 該從鍵盤讀的，還是得
\texttt{scanf}。

\keypoint{\texttt{printf} 用變數的值，\texttt{scanf} 用變數的位址。
\texttt{-Wall} 對漏寫 \texttt{\&}、以及使用未初始化變數，通常會先
喊一聲；那一聲值得當成錯誤來修。}

\begin{drills}[label=drills:basic-echo]
\item 為什麼要先宣告變數、再 \texttt{scanf}、最後才 \texttt{printf}，
順序顛倒會怎樣？
\item 範例 \thechapter.14（\texttt{sum-two.c}）輸入 \texttt{3 7}，會
印出什麼？
\item \texttt{scanf("\%d", n)}（漏寫 \texttt{\&}）編譯時，\texttt{gcc
-Wall} 通常會給什麼樣的警告？
\end{drills}

\section{讀取多個值與調整輸出格式}
\label{sec:basic-multi-read}

題目有時要求分多行讀入，或在同一行讀入多個數字但輸出順序不同。 兩種
常見片語如下：

\begin{code}{c}{同一行讀入兩個整數}
scanf("%d %d", &a, &b);
\end{code}

\begin{code}{c}{分三次各讀一個整數}
scanf("%d", &a);
scanf("%d", &b);
scanf("%d", &c);
\end{code}

同一行的 \texttt{"\%d \%d"} 會跳過中間的空白字元；分三次 \texttt{scanf}
則是每次讀到下一個整數為止。 兩者都能處理一行多個數字或每行一
個數字的輸入，但程式可讀性上，若題目明確說一行兩個整數，用一
次 \texttt{scanf} 讀齊通常比較清楚。

範例 \thechapter.17 讀入兩個整數，但輸出時先印第二個、再印第一個。 
只需要交換 \texttt{printf} 裡的變數順序，不需要真的交換變數的值。

\begin{example}{c}{(swap-two.c) 讀入兩個整數並以相反順序輸出}
#include <stdio.h>

int main(void) {
    int a, b;
    scanf("%d %d", &a, &b);

    printf("%d %d\n", b, a);
}
\end{example}

\iosample{
3 7 \\
}{
7 3 \\
}

範例 \thechapter.18 則示範分三次讀入三個整數：每次 \texttt{scanf} 各處
理一個 \texttt{\%d}，適合題目把三個數字分成三行給定的情況。

\begin{example}{c}{(read-three-sum.c) 分三次讀入三個整數}
#include <stdio.h>

int main(void) {
    int a, b, c;
    scanf("%d", &a);
    scanf("%d", &b);
    scanf("%d", &c);

    printf("%d\n", a + b + c);
}
\end{example}

\iosample{
1 \\
2 \\
3 \\
}{
6 \\
}

\stylepoint{輸出多個以空白分隔的整數時，格式字串寫成
\texttt{"\%d \%d\textbackslash n"}；最後一個數字後面要不要空白，要
對照題目範例輸出，不要自己多加。}

\begin{drills}[label=drills:basic-multi-read]
\item 用一次 \texttt{scanf("\%d \%d", \&a, \&b)} 跟分兩次
\texttt{scanf("\%d", \&a); scanf("\%d", \&b);}，讀到的結果一樣嗎？
\item 範例 \thechapter.17（\texttt{swap-two.c}）輸入 \texttt{3 7}，
輸出是什麼？靠的是交換變數的值，還是交換 \texttt{printf} 裡的順序？
\item 若題目把三個整數分三行給，該用一次 \texttt{scanf} 還是三次？
\end{drills}

\section{printf 的寬度與補零}
\label{sec:basic-width-padding}

除了 \texttt{\%d}，\texttt{printf} 還可以在 \texttt{\%} 與轉換字元之間
插入\emph{寬度}，控制至少印出幾個字元；寬度前面再加 \texttt{0}，不足
的部分會用前導零補齊，而不是空白。

\begin{code}{c}{將整數以至少 5 位數輸出，不足補零}
printf("%05d\n", id);
\end{code}

\texttt{\%05d} 裡的 \texttt{5} 是最小欄寬，\texttt{0} 表示用零補位。
例如 \texttt{id = 42} 只佔兩位，輸出會變成 \texttt{00042}。 若數字本
身已超過五位（例如 \texttt{123456}），不會被截斷，會完整印出六位。

範例 \thechapter.20 讀入一個學號或編號，依題目要求補零成五位。

\begin{example}{c}{(padded-id.c) 將整數補零成五位輸出}
#include <stdio.h>

int main(void) {
    int id;
    scanf("%d", &id);

    printf("%05d\n", id);
}
\end{example}

\iosample{
42 \\
}{
00042 \\
}

這種格式控制在處理固定長度的數字編碼時很常見。 例如日期裡的月
分、流水號。 算式怎麼從整數裡拆出各位數字，則屬於下一章的除法與餘數
技巧；這裡只處理已知一個整數，要怎麼漂亮地印出來。

\keypoint{\texttt{\%05d} 的 \texttt{0} 是補零旗標，\texttt{5} 是
最小寬度；兩者缺一不可，\texttt{\%5d} 只會用空白補位。}

\begin{drills}[label=drills:basic-width-padding]
\item \texttt{printf("\%05d\textbackslash n", 42);} 會印出什麼？
\item 若 \texttt{id = 123456}（超過 5 位數），\texttt{\%05d} 會把它
截斷成 5 位嗎？
\item \texttt{\%5d}（沒有前導 \texttt{0}）跟 \texttt{\%05d}，不足位數
時分別用什麼字元補齊？
\end{drills}

\section{double 型別與浮點輸入輸出}
\label{sec:basic-double}

有些輸出需要小數點（例如金額），這時候整數型別不夠用，要改用
\texttt{double}（雙精度浮點數）。 讀寫 \texttt{double} 時，格式字串
跟 \texttt{int} 不同：

\begin{code}{c}{讀入 double 並印到小數點後兩位}
double x;
scanf("%lf", &x);
printf("%.2f\n", x);
\end{code}

\texttt{scanf} 讀 \texttt{double} 必須用 \texttt{\%lf}（這裡的
\texttt{l} 代表 \texttt{long double} 家族，對 \texttt{double} 就是正
確的格式）；\texttt{printf} 印 \texttt{double} 則用 \texttt{\%f} 就
夠。 這是 C 語言留下來的不對稱寫法，之後浮點數那一章還會再從記憶體
表示的角度細講。

範例 \thechapter.22 讀入一個浮點數，四捨五入到小數點後兩位輸出。 這
裡的重點是 \texttt{\%lf}/\texttt{\%.2f} 的配對；至於整數分鐘換算
成含小數的元這類題目，分子該除以 \texttt{100.0} 還是 \texttt{100}，
則留到下一章討論型別轉換。

\begin{example}{c}{(read-double-print.c) 讀入浮點數並印出兩位小數}
#include <stdio.h>

int main(void) {
    double price;
    scanf("%lf", &price);

    printf("%.2f\n", price);
}
\end{example}

\iosample{
3.14159 \\
}{
3.14 \\
}

\texttt{\%.2f} 裡的 \texttt{.2} 表示小數位數；編譯器會做四捨五入。 若
不指定小數位數，\texttt{\%f} 預設常印出六位小數，對很多題目來說既冗
長又可能暴露浮點誤差。 養成依題目要求寫 \texttt{\%.Nf} 的習慣，可以少
掉不少格式不符的錯誤。

\stylepoint{讀 \texttt{double} 用 \texttt{\%lf}，印 \texttt{double} 用
\texttt{\%f} 或 \texttt{\%.Nf}；不要把 \texttt{scanf} 的格式直接套到
\texttt{printf} 上。}

C 對標點幾乎沒有容錯：少一個分號、括號不成對、字串少關一個直引號，
編譯器就無法確定你想表達什麼。 上一章的 \texttt{Hello, world!} 已經
示範過分號；這一章開始，每一個 \texttt{\%}、\texttt{\&}、逗號、雙引
號都各自有職務，不要把它們當成裝飾。 下一章會把加、減、乘、除、以
及整數除法的截斷，接到這個骨架上。 輸入輸出的管子先通，算式才有
地方可放。

\begin{drills}[label=drills:basic-double]
\item \texttt{scanf} 讀 \texttt{double} 用什麼格式？\texttt{printf}
印 \texttt{double} 用什麼格式？兩者一樣嗎？
\item \texttt{printf("\%.2f\textbackslash n", 3.14159);} 會印出什麼？
\item 若不指定小數位數，直接用 \texttt{\%f}，通常會印出幾位小數？
\end{drills}

\addcontentsline{toc}{section}{練習解答}
\section*{練習解答}

各節練習的答案收在這裡。 短題只給能一句話核對的結果。

\paragraph{第~\ref{sec:basic-printf}~節}
\begin{enumerate}
\item \texttt{1729}：兩次輸出直接黏在一起，中間沒有空白也沒有換行。
\item \texttt{"17 29\textbackslash n"}——空白要自己放進格式字串裡，
編譯器不會自動加。
\item \texttt{\textbackslash n}（換行字元）。
\end{enumerate}

\paragraph{第~\ref{sec:basic-int-vars}~節}
\begin{enumerate}
\item 型別、名字、值、位址；\texttt{scanf} 需要的是位址，所以要寫
\texttt{\&x}。
\item 內容不確定（殘留值）；不該直接拿去 \texttt{printf}，結果沒有保
證。
\item 不是同一個變數。 C 的識別字大小寫視為不同。
\end{enumerate}

\paragraph{第~\ref{sec:basic-init-assign}~節}
\begin{enumerate}
\item 一樣，最後 \texttt{n} 都是 \texttt{42}。
\item 不安全。 未初始化、未賦值的變數內容是不確定的殘留值，不該直接
使用。
\item \texttt{int a = 1, b = 2, c = 3;}
\end{enumerate}

\paragraph{第~\ref{sec:basic-scanf}~節}
\begin{enumerate}
\item \texttt{scanf("\%d", \&n);}
\item 因為 \texttt{scanf} 格式字串裡的空白（含 \texttt{\textbackslash
n}）會一直跳過輸入裡的空白字元，直到遇到下一個非空白字元；多打的
\texttt{\textbackslash n} 會讓程式等使用者多輸入一次才繼續。
\item 因為 \texttt{scanf} 需要知道要把讀到的值\emph{寫到哪裡}（位
址），只給變數名（值）它無從得知該寫進哪一格記憶體。
\end{enumerate}

\paragraph{第~\ref{sec:basic-echo}~節}
\begin{enumerate}
\item 因為要先有變數才能讀入，要先讀入才有值可印；顛倒的話，
\texttt{printf} 會印出尚未寫入的內容，或 \texttt{scanf} 沒有地方可
存。
\item \texttt{10}（\texttt{3 + 7}）。
\item 格式期待 \texttt{int *}（位址）卻收到 \texttt{int}（值），並提
醒 \texttt{n} 在呼叫時尚未初始化。
\end{enumerate}

\paragraph{第~\ref{sec:basic-multi-read}~節}
\begin{enumerate}
\item 一樣：兩種寫法都能把兩個以空白（或換行）分隔的整數分別讀進
\texttt{a}、\texttt{b}。
\item 輸出 \texttt{7 3}；靠的是交換 \texttt{printf} 裡兩個變數的順
序，值本身沒有互換。
\item 分三次 \texttt{scanf("\%d", ...)}，每次讀一個。
\end{enumerate}

\paragraph{第~\ref{sec:basic-width-padding}~節}
\begin{enumerate}
\item \texttt{00042}。
\item 不會截斷，會完整印出 \texttt{123456}（六位）。
\item \texttt{\%5d} 用空白補位；\texttt{\%05d} 用 \texttt{0} 補位。
\end{enumerate}

\paragraph{第~\ref{sec:basic-double}~節}
\begin{enumerate}
\item 不一樣：\texttt{scanf} 讀 \texttt{double} 要用 \texttt{\%lf}，
\texttt{printf} 印則用 \texttt{\%f}（或 \texttt{\%.Nf}）。
\item \texttt{3.14}。
\item 預設通常印出六位小數。
\end{enumerate}
